\section{Загальний опис системи}\label{sec:opis}

\subsection{Архітектура системи та перспектива продукту}\label{sec:opis-arkhitektura}
Система проєктується як сучасний вебзастосунок із класичною клієнт-серверною архітектурою~\sked{K9}. Клієнтська частина функціонує у веббраузері користувача та забезпечує відображення афіші, розкладу сеансів, схеми залу та інтерфейсу особистого кабінету~\sked{K9}. Серверна частина надає програмний інтерфейс (API), виконує бізнес-логіку та забезпечує транзакційну обробку операцій~\sked{K9}.

Ключовим архітектурним принципом є безумовна авторитетність серверної частини: перевірка прав доступу, коректності вхідних параметрів та бізнес-правил бронювання здійснюється на сервері незалежно від валідацій на стороні вебклієнта~\sked{K40}. Усі дані користувачів, фільмів, кінозалів, розкладу сеансів та створених бронювань зберігаються персистентно й гарантовано відновлюються після завершення роботи або перезапуску серверних процесів~\sked{K41}.

\begin{figure}[htbp]
\centering
\makebox[\textwidth][c]{%
\begin{tikzpicture}[node distance=10mm and 32mm]
  \node[sked box] (client) {Вебклієнт\\(Браузер / SPA)};
  \node[sked box, right=of client] (server) {Серверний застосунок\\(API \& Бізнес-правила)};
  \node[sked box, right=of server] (db) {Персистентне сховище\\(База даних / Сховище)};
  \draw[sked arrow] ([yshift=2.5mm]client.east) -- node[sked label, above=1pt] {Запити API} ([yshift=2.5mm]server.west);
  \draw[sked arrow] ([yshift=-2.5mm]server.west) -- node[sked label, below=1pt] {Відповіді \& Статуси} ([yshift=-2.5mm]client.east);
  \draw[sked arrow] ([yshift=2.5mm]server.east) -- node[sked label, above=1pt] {Атомарні транзакції} ([yshift=2.5mm]db.west);
  \draw[sked arrow] ([yshift=-2.5mm]db.west) -- node[sked label, below=1pt] {Персистентні дані} ([yshift=-2.5mm]server.east);
\end{tikzpicture}%
}
\caption{Контекстна клієнт-серверна архітектура системи}\label{fig:architecture}
\end{figure}

\subsection{Класи користувачів та характеристики ролей}\label{sec:opis-korystuvachi}
Система чітко розмежовує два класи користувачів на основі моделі рольового доступу (RBAC): Користувач (зареєстрований клієнт кінотеатру) та Адміністратор~\sked{K1}. 

\begin{itemize}
    \item \textbf{Користувач (клієнт):} зареєстрований відвідувач, що автентифікувався за email та паролем і має унікальний внутрішній числовий/ідентифікаційний ID. Має право переглядати репертуар, створювати бронювання на своє ім'я та переглядати чи скасовувати виключно власні бронювання~\sked{K1}. Звичайний користувач суворо позбавлений доступу до адміністративних операцій і не має технічної можливості призначити собі роль адміністратора~\sked{K27}.
    \item \textbf{Адміністратор:} співробітник кінотеатру, який працює в тій самій системі через спеціальний відокремлений розділ керування~\sked{K26}. Відповідає за наповнення довідників фільмів, залів і сеансів, а також за обробку позаштатних ситуацій (скасування сеансів)~\sked{K26}. Для першої версії передбачено створення рівно одного адміністративного облікового запису під час початкової конфігурації системи; підтримка керування кількома адміністраторами винесена за межі v1~\sked{K28}.
\end{itemize}

У \tabref{tab:rbac-matrix} наведено матрицю доступу до операцій системи для кожної ролі.

\begin{table}[htbp]
\caption{Матриця прав доступу ролей (RBAC)}\label{tab:rbac-matrix}
\centering
\begin{tabularx}{\textwidth}{@{}X >{\centering\arraybackslash}p{2.8cm} >{\centering\arraybackslash}p{2.6cm} >{\centering\arraybackslash}p{2.8cm}@{}}
\toprule
\textbf{Операція / Ресурс} & \textbf{Анонімний відвідувач} & \textbf{Користувач (Клієнт)} & \textbf{Адміністратор} \\
\midrule
Перегляд фільмів та сеансів & Дозволено & Дозволено & Дозволено \\
Перегляд схеми залу & Дозволено & Дозволено & Дозволено \\
Реєстрація та вхід (email+пароль) & Дозволено & Вже виконано & Вже виконано \\
Створення бронювання (1--6 місць) & Заборонено & Дозволено & Заборонено (у v1) \\
Перегляд власних бронювань & Заборонено & Дозволено & Дозволено \\
Скасування власних бронювань & Заборонено & Дозволено & Не застосовується \\
Керування фільмами, залами, сеансами & Заборонено & Заборонено & Дозволено \\
Явне скасування сеансу з причиною & Заборонено & Заборонено & Дозволено \\
Перегляд операційних логів & Заборонено & Заборонено & Дозволено \\
\bottomrule
\end{tabularx}
\end{table}

\subsection{Середовище функціонування та припущення}\label{sec:opis-seredovyshche}
Клієнтська частина повинна стабільно функціонувати в усіх сучасних веббраузерах (Google Chrome, Mozilla Firefox, Safari, Microsoft Edge). Серверна частина розгортається в середовищі, що підтримує роботу API та зв'язок зі сховищем даних~\sked{K9}.

Ключовим бізнес-припущенням є те, що бронювання слугує попереднім резервуванням місць, а безпосередній розрахунок грошима за квитки та отримання фіскальних чеків відбуваються в касі кінотеатру перед початком сеансу поза межами автоматизованої системи~\sked{K35}.

\subsection{Загальні обмеження проєктування}\label{sec:opis-obmezhennia}
При проєктуванні системи встановлюються такі засадничі обмеження:
\begin{itemize}
    \item \textbf{Ліміт бронювання:} одне бронювання повинно містити від 1 до 6 місць включно (максимально допустима кількість --- 6 місць)~\sked{K4}.
    \item \textbf{Відсутність утримання місць:} вибір місць на схемі залу у вебінтерфейсі не створює тимчасового блокування (hold) і не гарантує доступності місць до завершення атомарної операції запису~\sked{K31}.
    \item \textbf{Відмова від таймаутів у v1:} у версії 1 механізм блокування місць за таймаутом не реалізується; натомість застосовується перевірка конфліктів у момент оформлення~\sked{K32}.
\end{itemize}
